% ==========================================================================
% GENERATED by scripts/split_paper.py. Do not edit this file.
%
% Source: imported_papers/probe-norm/content/body.tex, content/supplement.tex
% Part 10: Code for the Protocol
%
% Edits belong upstream, in the published paper. To change what the
% chapter includes, change imports/bounded.toml and re-run
% the splitter.
% ==========================================================================
\section{Code for the Protocol}
\label{app:code}
The module \texttt{probenorm.protocol}, released with this paper, runs the four steps of Section~\ref{sec:decision} for one model or a pair of models, from the command line or from Python. The listing shows these steps with the functions of the \texttt{probenorm} package. It assumes reference encodings \texttt{Z} with states \texttt{Y} and episode ids \texttt{groups}, predictions \texttt{P} at one horizon with their true states \texttt{y}, the initial encodings \texttt{Z0}, and the encoded targets \texttt{Zt}. Null rollouts and the start-encoding probe depend on the model and the data, so the listing takes them as given.

\begin{tcolorbox}[breakable,colback=codebg,colframe=codebg,boxrule=0pt,arc=2pt,left=4pt,right=4pt,top=4pt,bottom=4pt]
\begin{lstlisting}
import numpy as np
from probenorm.probes import fit_probe, direction, mse
from probenorm.reproject import ReferenceBank, rescale
from probenorm.checks import radius_check, support_check, random_directions

# 1. Check the probe on encodings.
fits = [fit_probe(Z, Y, groups, seed=s) for s in range(3)]
probes = [probe for probe, _ in fits]
r2 = np.mean([info["r2"] for _, info in fits])   # low: do not apply the probe to predictions
rc = radius_check(Z, Y, groups)
gain = np.mean(rc["aug_std"]["r2"]) - np.mean(rc["dir_std"]["r2"])

# 2. Compare raw and re-projected errors.
def error(V):
    return np.mean([mse(probe(V), y) for probe in probes])

bank = ReferenceBank(Z)
rho = bank.radius(P)                 # or the analytic radius of the encoder's normalization
raw, rep = error(P), error(rescale(P, rho))
factor = rep / raw
read_rep = gain <= 0.05 and factor < 0.9    # otherwise keep the raw error

# 3. Test against the null controls: the model is rejected unless every comparator is worse;
#    a ranking is rejected when both models are, or when it puts a rejected one first.
R = random_directions(P)
comparators = {"persistence": error(Z0),
               "reference mean": mse(np.broadcast_to(Y.mean(0), y.shape), y),
               "start-encoding probe": mse(start_probe(Z0), y),
               "random directions": error(rescale(R, bank.radius(R))),
               **{name: error(rescale(V, bank.radius(V)))
                  for name, V in null_rollouts.items()}}
rejected = rep >= min(comparators.values())

# 4. Check support and the other probes: inconclusive below one half supported, and for a
#    ranking unless the re-projected and normalized probes agree and no other probe
#    reverses them (repeat per model).
supported = support_check(bank, P, Zt)[0].mean()
normalized, _ = fit_probe(Z, Y, groups, feature=direction)
scores = {"re-projected": rep, "normalized": mse(normalized(P), y)}
\end{lstlisting}
\end{tcolorbox}
